\documentclass[serif, xcolor={svgnames, table}, usepdftitle=false]{beamer}

% Beamer
\usetheme{CambridgeUS}
\usecolortheme{beaver}

% Encoding
\usepackage[utf8]{inputenc}
\usepackage[T2A]{fontenc}

% Languages and bibliographies
\usepackage[macedonian]{babel}
\usepackage{csquotes}
\usepackage[hyperref, backref, backend=biber]{biblatex}

% Mathematics
\usepackage{mathtools}
\usepackage{amsthm}
\usepackage{amssymb}

% Pseudocode
\usepackage{algorithm}
\usepackage{algpseudocode}
\usepackage{varwidth}

% Microtype
\usepackage[babel]{microtype}

% Graphics and captioning
\usepackage{graphicx}
\usepackage{caption}
\usepackage{subcaption}

% Tables
\usepackage{booktabs}

% Numbers
\usepackage{siunitx}
\usepackage{nth}

% TikZ
\usepackage{tikz}
\usepackage{pgfplots}

%
% Package setup and custom commands
%

\sisetup{%
  detect-all,
  detect-display-math,
  group-separator=\,,
  output-decimal-marker={,}
}

\mathtoolsset{mathic=true}

\newcommand*{\RR}{\ensuremath{\mathbb{R}}}

\DeclarePairedDelimiter{\ceil}{\lceil}{\rceil}
\DeclarePairedDelimiter{\floor}{\lfloor}{\rfloor}
\DeclarePairedDelimiter{\round}{\lceil}{\rfloor}
\DeclarePairedDelimiter{\dotp}{\langle}{\rangle}
\DeclarePairedDelimiter{\abs}{\lvert}{\rvert}
\DeclarePairedDelimiter{\norm}{\lVert}{\rVert}

\usetikzlibrary{shapes, arrows}
\pgfplotsset{compat=1.13}

\graphicspath{{Figures/}{../}}

%
% Localization
%

\deftranslation[to=macedonian]{Theorem}{Теорема}
\deftranslation[to=macedonian]{theorem}{теорема}
\deftranslation[to=macedonian]{Definition}{Дефиниција}
\deftranslation[to=macedonian]{definition}{дефиниција}

%
% Title
%

\title{Хаусхолдеров метод}
\subtitle{Алгоритам за наоѓање корени на функции}
\author[Дарио Ѓорѓевски]{%
  Дарио Ѓорѓевски\inst{1}\\%
  \href{mailto:gjorgjevski.dario@students.finki.ukim.mk}%
  {\texttt{gjorgjevski.dario@students.finki.ukim.mk}}
}
\institute[ФИНКИ]{%
  \inst{1}Факултет за компјутерски науки и инженерство\\%
  Универзитет Св.\ Кирил и Методиј, Скопје
}
\date{\today}

\hypersetup{%
  unicode,
  pdfinfo={%
    Title={Хаусхолдеров метод},
    Author={Дарио Ѓорѓевски},
    Subject={Нумеричка математика},    
    Keywords={Математика, Нумеричка математика, Нумеричка анализа}
  }
}

%
% Logo
%

\logo{\includegraphics[height=0.66cm]{Logo.png}}

%
% Table of contents
%

\AtBeginSection[]{%
  \begin{frame}{Содржина}
    \tableofcontents[currentsection, hideothersubsections]
  \end{frame}
}

\begin{document}

\begin{frame}
  \titlepage
\end{frame}

\section{Вовед и опис на методот}

\begin{frame}{Мотивација}
  \only<1>{%
    \framesubtitle{Апроксимација на \((1 / {f})(x)\) со геометриски ред}
    
    Нека \(f(x) \colon \RR \to \RR\) е непрекинато диференцијабилна функција со
    корен со кратност \(1\) во \(x = a\).

    Со други зборови, \(f(a) = 0\) и \(f'(a) \ne 0\).

    Ова значи дека \((1 / {f})(x)\) има пол со кратност \(1\) во \(x = a\), и во
    околината на \(a\) се однесува како \(1 / {(x - a)}\).

    Приближно добиваме
    \begin{equation}\label{eq:approximation}
      \begin{aligned}
        \frac{1}{f(x)} = \frac{1}{f(x) - f(a)}
        &= \frac{x - a}{f(x) - f(a)} \frac{-1}{a (1 - x / {a})} \\
        &\approx \frac{-1}{a f'(a)} \sum\limits_{k = 0}^{\infty}
        \left(\frac{x}{a}\right)^k\text{.}
      \end{aligned}
    \end{equation}
  }

  \only<2>{%
    \framesubtitle{Наоѓање на \(a\)}

    Коефициентот пред \(x^d\) во~\eqref{eq:approximation} е \(C a^{-d}\)
    \(\Rightarrow\) вредноста на \(a\) можеме да ја добиеме така што
    коефициентот пред \(x^{d - 1}\) ќе го поделиме со коефициентот пред \(x^d\).

    \begin{alertblock}{Клучна идеја}
      Геометрискиот ред од~\eqref{eq:approximation} дава апроксимација на
      Тајлоровата експанзија на \((1 / {f})(x)\) во околината на \(a\).  Ова ни
      овозможува да го апроксимираме \(a\) со делење на соодветните коефициенти
      во Тајлоровата експанзија на \((1 / {f})(b + x)\):
      \begin{equation*}
        a \approx b + \frac{(1 / {f})^{(d - 1)}(b)}{(d - 1)!} \frac{d!}{(1 / {f})^{(d)}(b)}
        = b + d \frac{(1 / {f})^{(d - 1)}(b)}{(1 / {f})^{(d)}(b)}\text{.}
      \end{equation*}
    \end{alertblock}
  }
\end{frame}

\begin{frame}{Опис на Хаусхолдеровиот метод}
  \begin{definition}[Хаусхолдеров метод од ред \(d\)]
    Нека е дадена непрекинато диференцијабилна функција
    \(f(x) \colon \RR \to \RR\).  За наоѓање на корен на \(f\), вршиме низа
    итерации
    \begin{equation}\label{eq:householder-method}
      x_{n + 1} = x_n + d \frac{(1 / {f})^{(d - 1)}(x_n)}{(1 / {f})^{(d)}(x_n)}\text{,}
    \end{equation}
    почнувајќи од претпоставка \(x_0\).
  \end{definition}

  Доколку \(f\) е \((d + 1)\) пат непрекинато диференцијабилна и \(a\) е корен
  со кратност \(1\), тогаш
  \[
    \abs{x_{n + 1} - a} \le K \abs{x_n - a}^{d + 1}\text{, за некое }K > 0\text{.}
  \]
\end{frame}

\begin{frame}{Врска со други методи}
  Хаусхолдеровиот метод е генерален; други методи се негови специјални случаи.
  \begin{itemize}
  \item Кога \(d = 1\) го добиваме Њутн--Рафсоновиот метод:
    \begin{align*}
      x_{n + 1} &= x_n + \frac{(1 / {f})(x_n)}{(1 / {f})'(x_n)} \\
                &= x_n + \frac{1}{f(x_n)} \left(\frac{-f'(x_n)}{f(x_n)^2}\right)^{-1} \\
                &= x_n - \frac{f(x_n)}{f'(x_n)}\text{.}
    \end{align*}
  \item Кога \(d = 2\) го добиваме Халеевиот метод.
  \end{itemize}
\end{frame}

\section{Пример на употреба}

\begin{frame}{Решавање на \(y^3 - 2 y - 5 = 0\)}
  \begin{exampleblock}{Употреба на Хаусхолдеровиот метод}
    Да претпоставиме дека сакаме да најдеме нула на равенката
    \(y^3 - 2 y - 5 = 0\).  Графички можеме да забележиме дека една нула се
    наоѓа во околината на \(y = 2\).  Оттука, земајки \(y \coloneqq x + 2\),
    равенката станува
    \begin{equation}\label{eq:solve-for}
      0 = f(x) = -1 + 10 x + 6 x^2 + x^3\text{.}
    \end{equation}
  \end{exampleblock}

  Следен чекор е наоѓање на Тајлорова експанзија на \((1 / {f})(x)\) околу
  \(0\), каде што \(f(x)\) е дадена со~\eqref{eq:solve-for}.
\end{frame}

\begin{frame}{График на \(y^3 - 2 y - 5\)}
  \begin{figure}
    \centering
    \begin{tikzpicture}
      \begin{axis}[ 
        xlabel=\(y\),
        axis lines=middle
        ]
        \addplot[blue, no marks]{x^3 - 2*x - 5};
      \end{axis}
    \end{tikzpicture}
  \end{figure}
\end{frame}

\begin{frame}{Тајлорова експанзија на \((1 / {f})(x)\)}
  Лесно е да се пресмета дека ако \(f(x) = -1 + 10 x + 6 x^2 + x^3\), тогаш
  Тајлоровата експанзија на \((1 / {f})(x)\) околу \(0\) е
  \begin{equation*}
    \begin{aligned}
      (1 / {f})(x) = \, & {-}\num{1} - \num{10} x - \num{106} x^2 - \num{1121} x^3 - \num{11856} x^4 \\
      & {-} \num{125392} x^5 - \num{1326177} x^6 - \num{14025978} x^7 - \num{148342234} x^8 \\
      & {-}\num{1568904385} x^9 - \num{16593123232} x^{10} + \mathcal{O}(x^{11})\text{.}
    \end{aligned}
  \end{equation*}

  Примената на Хаусхолдеровиот метод од различен ред во \(x = 0\) се сведува на
  делење соседни коефициенти во Тајлоровата експанзија.
\end{frame}

\begin{frame}{Вредности на \(x_1\) за различни редови на методот}
  \begin{table}
    \centering
    \begin{tabular}{S[table-format=2.0] S[table-format=1.33]}
      \toprule
      {\(d\)} & {\(x_1\)} \\
      \midrule
      1 & 0.100000000000000000000000000000000 \\
      2 & 0.094339622641509433962264150943396 \\
      3 & 0.094558429973238180196253345227476 \\
      4 & 0.094551282051282051282051282051282 \\
      5 & 0.094551486538216154140615031261963 \\
      6 & 0.094551481438752142436492263099119 \\
      7 & 0.094551481543746895938379484125813 \\
      8 & 0.094551481542336756233561913325371 \\
      9 & 0.094551481542324837086869382419375 \\
      10 & 0.094551481542326678478801765822985 \\
      \bottomrule
    \end{tabular}
  \end{table}
\end{frame}

\begin{frame}{Експлицитни формули за \(x_{n + 1}\) за пониски редови}
  Доколку наоѓаме изводи на \(f(x)\), можеме да видиме дека
  \begin{align*}
    f'(x) &= 10 + 12 x + 3 x^2 \\
    f''(x) &= 12 + 6 x \\
    f'''(x) &= 6\text{.}
  \end{align*}

  Од друга страна, може да се изведат експлицитни формули за \(x_{n + 1}\) за
  пониски редови:
  \begin{itemize}
  \item Кога \(d = 1\), \(x_{n + 1} = x_n - f / {f'}\).
  \item Кога \(d = 2\), \(x_{n + 1} = x_n + (-2 f f') / {(2 f'^2 - f f'')}\).
  \item Кога \(d = 3\),
    \(\displaystyle x_{n + 1} = x_n - \frac{6 f f'^2 - 3 f^2 f''}{6 f'^3 - 6 f
      f' f'' + f^2 f'''}\).
  \end{itemize}
\end{frame}

\begin{frame}{Вредности на \(x_2, x_3, x_4, x_5, x_6\) за \(d = 1\)}
  \begin{table}
    \centering
    \newcolumntype{C}{>{$}c<{$}}
    \begin{tabular}{C S[table-format=1.34]}
      \toprule
      x_n & {Вредност} \\
      \midrule
      x_1 & 0.1 \\
      x_2 & 0.09456812110419 \\
      x_3 & 0.09455148169819930297 \\
      x_4 & 0.09455148154232659149606485 \\
      x_5 & 0.09455148154232659148238654057931 \\
      x_6 & 0.09455148154232659148238654057931 \\
      \bottomrule
    \end{tabular}
  \end{table}
\end{frame}

\begin{frame}{Вредности на \(x_2, x_3, x_4, x_5\) за \(d = 2\)}
  \begin{table}
    \centering
    \newcolumntype{C}{>{$}c<{$}}
    \begin{tabular}{C S[table-format=1.31]}
      \toprule
      x_n & {Вредност} \\
      \midrule
      x_1 & 0.09433962264151 \\
      x_2 & 0.09455148154016421472 \\
      x_3 & 0.09455148154232659148238654 \\
      x_4 & 0.0945514815423265914823865405793 \\
      x_5 & 0.0945514815423265914823865405793 \\
      \bottomrule
    \end{tabular}
  \end{table}
\end{frame}

\begin{frame}{Вредности на \(x_2, x_3, x_4\) за \(d = 3\)}
  \begin{table}
    \centering
    \newcolumntype{C}{>{$}c<{$}}
    \begin{tabular}{C S[table-format=1.32]}
      \toprule
      x_n & {Вредност} \\
      \midrule
      x_1 & 0.09455842997324 \\
      x_2 & 0.094551481542326591482567 \\
      x_3 & 0.09455148154232659148238654057931 \\
      x_4 & 0.09455148154232659148238654057931 \\
      \bottomrule
    \end{tabular}
  \end{table}
\end{frame}

\section{Формално изведување}

\begin{frame}{Падеева апроксимација}
  \begin{definition}[Падеева апроксимација]
    За дадена функција \(f\) и два цели броја \(m \ge 0\) и \(n \ge 1\),
    \emph{Падеева апроксимација} од ред \([m / {n}]\) е рационалната функција
    \[
      R(x) = \frac{\sum\nolimits_{i = 0}^{m} a_i x^i}{1 + \sum\nolimits_{j =
          1}^{n} b_j x^j}\text{,}
    \]
    која со \(f\) „се согласува“ до највисок можен ред, односно
    \(f^{(k)}(0) = R^{(k)}(0)\) за сите \(0 \le k \le m + n\).
  \end{definition}
  Bottom line: апроксимацијата може да се најде преку решавање систем равенки
  добиен од Тајлоровата експанзија на \(f\) околу \(0\).
\end{frame}

\begin{frame}{Дефинирање на итерацијата}
  За да го изведеме Хаусхолдеровиот метод од ред \(d\), ќе започнеме со Падеева
  апроксимација на \(f(x + h)\) од ред \([1 / {(d - 1)}]\).  Таа има облик
  \begin{equation}\label{eq:pade-approximation}
    f(x + h) = \frac{a_0 + h}{b_0 + b_1 h + \ldots + b_{d - 1} h^{d - 1}} +
    \mathcal{O}(h^{d + 1})\text{.}
  \end{equation}
  
  Имено, знаеме дека \(b_d = 0\).  Итерацијата на Хаусхолдеровиот метод се
  добива преку наоѓање на нулата во \(h = -a_0\).
\end{frame}

\begin{frame}{Наоѓање на Падеевата апроксимација}
  Најлесен начин за наоѓање на потребната Падеева апроксимација е преку
  Тајлорова експанзија на \((1 / {f})(x + h)\) околу \(0\):
  \begin{equation}\label{eq:taylor-expansion}
    \begin{aligned}
      (1 / {f})(x + h) &= (1 / {f})(x) + (1 / {f})'(x) h + \ldots \\
      &+ (1 / {f})^{(d - 1)}(x) \frac{h^{d - 1}}{(d - 1)!} + (1 / {f})^{(d)}(x)
      \frac{h^d}{d!}\text{.}
    \end{aligned}
  \end{equation}

  Можеме да воочиме дека~\alert{\eqref{eq:taylor-expansion} треба да биде
    еднакво на реципрочната вредност на~\eqref{eq:pade-approximation}}.
\end{frame}

\begin{frame}{Наоѓање на нулата во \(h = -a_0\)}
  Користејќи ги својствата на Падеевата апроксимација, добиваме
  \[
    0 = b_d = a_0 (1 / {f})^{(d)}(x) \frac{1}{d!} + (1 / {f})^{(d - 1)}(x)
    \frac{1}{(d - 1)!}\text{.}
  \]

  Оттука веднаш се добива
  \[
    h = -a_0 = d \frac{(1 / {f})^{(d - 1)}(x)}{(1 / {f})^{(d)}(x)}\text{,}
  \]
  што ја дава итерацијата во Хаусхолдеровиот метод.
\end{frame}

\begin{frame}[plain, c]
\begin{center}
  {\huge КРАЈ}\\
  Прашања?
\end{center}
\end{frame}

\end{document}

%%% Local Variables:
%%% mode: latex
%%% fill-column: 80
%%% TeX-command-extra-options: "-shell-escape"
%%% TeX-engine: default
%%% TeX-master: t
%%% End: